% Problem record op_ddfbbe7ca0ead0c8. This TeX file is derived from
% database/problems_json/op_ddfbbe7ca0ead0c8.json by scripts/sync-tex.mjs; edit the JSON record
% and rerun the script rather than editing this file.
\section{Non-additive ((6,2,3)) qubit code with transversal T}
\label{sec:op_ddfbbe7ca0ead0c8}

\paragraph{Problem.}
Does there exist a genuinely non-additive two-dimensional six-qubit code of
distance three on which a tensor product of single-qubit unitaries implements
the logical $T$ gate exactly?  Write $\mathcal P_n:=\{I,X,Y,Z\}^{\otimes n}$
for the phase-free Hermitian Pauli basis, where the weight
$\operatorname{wt}(E)$ of $E\in\mathcal P_n$ counts its nonidentity tensor
factors.  An $((n,K,d))_2$ code is a $K$-dimensional subspace of
$(\mathbb C^2)^{\otimes n}$ that detects every Pauli operator of weight less
than $d$; degenerate codes are allowed.  A $((6,2,3))_2$ code is thus the range
of an isometry $V:\mathbb C^2\to(\mathbb C^2)^{\otimes6}$, with code projector
$P:=VV^\dagger$, satisfying
\begin{equation}
  V^\dagger V=I_2,
  \qquad
  V^\dagger EV=c_EI_2
  \quad\text{for all }E\in\mathcal P_6\text{ with }\operatorname{wt}(E)<3,
  \qquad c_E\in\mathbb R.
  \label{eq:ddfb-kl}
\end{equation}
The code is \emph{genuinely non-additive} if no tensor product of single-qubit
unitaries, even combined with a permutation of the qubits, maps it onto a Pauli
stabilizer code, that is, onto the joint $+1$ eigenspace of an abelian subgroup
of the Pauli group.  The logical gate and its transversal realization are
\begin{equation}
  T:=\operatorname{diag}(1,e^{i\pi/4}),
  \qquad
  \Bigl(\bigotimes_{j=1}^{6}U_j\Bigr)V=e^{i\phi}VT,
  \qquad U_1,\ldots,U_6\in U(2),\quad \phi\in\mathbb R.
  \label{eq:ddfb-transversal}
\end{equation}
The local factors $U_j$ may differ from one another and need not be physical
$T$ gates.  No ancillary qubits, measurements, code switching, or additional
transversal logical gates are assumed.  The question asks whether some
isometry satisfying Eq.~\eqref{eq:ddfb-kl} has a genuinely non-additive range
and admits unitaries $U_1,\ldots,U_6$ and a phase $\phi$ satisfying
Eq.~\eqref{eq:ddfb-transversal}.

\subsection*{Status}

Unsolved

\subsection*{Source}

Derived from a documented gap.  Zhang, Wu, Huang, and Zeng realize
transversal $T$ on seven qubits (Section~6.3.1), and their Table~2 lists the
cyclic group $\mathrm C_{16}$, generated by the determinant-one representative
of $T$, among the transversal groups for which no $((6,2,3))$ code was found
\sourcecite{ref:ddfb-zwhz}{ZWHZ25}.  The source neither poses the six-qubit
question explicitly nor claims nonexistence; the requirement of genuine
non-additivity belongs to the present formulation.

\subsection*{Progress}

\begin{itemize}
  \item By the Knill--Laflamme conditions, Eq.~\eqref{eq:ddfb-kl} is equivalent to
  exact correction of the error set
  $\mathcal F:=\{I^{\otimes6}\}\cup\{E\in\mathcal P_6:\operatorname{wt}(E)=1\}$,
  \begin{equation}
    PF_a^\dagger F_bP=\Lambda_{ab}P
    \qquad(F_a,F_b\in\mathcal F),
    \label{eq:ddfb-correction}
  \end{equation}
  because the products $F_a^\dagger F_b$ span the operators of weight at most
  two.  The question therefore concerns arbitrary single-qubit errors, not a
  particular noise channel \sourcecite{ref:ddfb-kl}{KL97}.

  \item Du, Zhang, Poon, and Zeng introduced the signature norm of a code.  For the
  rank-two projector $P_n$ of an $((n,2,3))_2$ code its square is
  \begin{equation}
    s(P_n)^2:=\frac14
    \sum_{\substack{E\in\mathcal P_n\\ 1\leq\operatorname{wt}(E)\leq2}}
    \lvert\operatorname{Tr}(P_nE)\rvert^2 .
    \label{eq:ddfb-signature}
  \end{equation}
  They prove that it is invariant under local unitaries; it is also invariant
  under qubit permutations.  For a stabilizer code each coefficient
  $\operatorname{Tr}(P_nE)/2$ lies in $\{0,\pm1\}$, so $s(P_n)^2$ is an integer
  (Appendix~C of the arXiv version).  A noninteger value of
  Eq.~\eqref{eq:ddfb-signature} therefore certifies genuine non-additivity in
  the sense of the statement \sourcecite{ref:ddfb-dzpz}{DZPZ26}.

  \item Angl\`es Munn\'e and Huber give an exact rational semidefinite-programming
  certificate that every $((6,2,3))_2$ code is impure (Section~4.2, item~(e);
  Table~1 marks the optimal size $K=2$ at block length $6$ and distance $3$ as
  impure).  Equivalently, $s(P_6)^2>0$ for every isometry satisfying
  Eq.~\eqref{eq:ddfb-kl}, so some Pauli operator of weight one or two has a
  nonzero expectation value on any candidate code
  \sourcecite{ref:ddfb-amh}{AMH26}.

  \item Seven qubits suffice.  Section~6.3.1 of Zhang, Wu, Huang, and Zeng gives an
  explicit exact $((7,2,3))_2$ code with two free phases,
  $\lvert1_L\rangle=X^{\otimes7}\lvert0_L\rangle$, and angle vector
  $(1,2,2,2,2,3,3)$: the operator $X^{\otimes7}$ and a tensor product of
  single-qubit $Z$ rotations implement logical $X$ and, up to global phase,
  $T^\dagger=\operatorname{diag}(1,e^{-i\pi/4})$.  Hence $T=(T^\dagger)^7$ is
  transversal as well.  The example has
  \begin{equation}
    s(P_7)^2=\frac{21}{8},
    \label{eq:ddfb-seven}
  \end{equation}
  so Eq.~\eqref{eq:ddfb-seven} certifies that it is genuinely non-additive
  \sourcecite{ref:ddfb-zwhz}{ZWHZ25},
  \sourcecite{ref:ddfb-dzpz}{DZPZ26}.

  \item Six qubits already support a different non-Clifford phase.  Section~4.1 of
  the same paper gives an exact $((6,2,3))_2$ code with transversal logical
  $\operatorname{diag}(1,e^{-2\pi i/5})$, up to global phase, and signature norm
  $\sqrt{0.84}$, that is, $s(P_6)^2=21/25$; this code is therefore genuinely
  non-additive.  Its logical phase has projective order $5$, whereas $T$ has
  projective order $8$.  Table~2 lists $\mathrm C_{16}$, generated by
  $e^{-i\pi/8}T$, among the targets for which no $((6,2,3))$ code was found; the
  caption states only that no such code has been found so far
  \sourcecite{ref:ddfb-zwhz}{ZWHZ25}.

  \item He, Lu, and Zeng (arXiv version~3, April 2026) exactly resolve the twelve
  seven-qubit candidates that survive the subset-sum and linear-programming
  filters in the complementary binary-dihedral ansatz, in which the physical
  gates are diagonal phase gates with angles in multiples of $2\pi/8$ and
  $\lvert1_L\rangle=X^{\otimes7}\lvert0_L\rangle$:
  \begin{equation}
    12\ \text{filtered candidates}
    =10\ \text{realizable candidates}+2\ \text{excluded candidates}.
    \label{eq:ddfb-classification}
  \end{equation}
  Section~IV and Appendices~F--G give the exact constructions and no-go proofs,
  which the authors report as verified in Lean.  This restricted seven-qubit
  classification does not settle the unrestricted six-qubit question; the same
  paper's catalogue on at most six qubits concerns distance two
  \sourcecite{ref:ddfb-hlz}{HLZ25}.
\end{itemize}

\subsection*{References}

\begin{enumerate}[label={},leftmargin=4.8em,labelsep=0.5em]
  \footnotesize
  \item[\textup{[KL97]}]\label{ref:ddfb-kl}
  E. Knill and R. Laflamme, ``Theory of quantum error-correcting codes,'' \emph{Physical Review A} \textbf{55}, 900--911 (1997). \href{https://doi.org/10.1103/PhysRevA.55.900}{doi:10.1103/PhysRevA.55.900}; \href{https://arxiv.org/abs/quant-ph/9604034}{arXiv:quant-ph/9604034}.

  \item[\textup{[DZPZ26]}]\label{ref:ddfb-dzpz}
  M. Du, C. Zhang, Y. T. Poon, and B. Zeng, ``Characterizing quantum codes via the coefficients in Knill--Laflamme conditions,'' \emph{npj Quantum Information} \textbf{12}, 13 (2026). \href{https://doi.org/10.1038/s41534-025-01155-1}{doi:10.1038/s41534-025-01155-1}; \href{https://arxiv.org/abs/2410.07983}{arXiv:2410.07983}.

  \item[\textup{[ZWHZ25]}]\label{ref:ddfb-zwhz}
  C. Zhang, Z. Wu, S. Huang, and B. Zeng, ``Transversal Gates in Nonadditive Quantum Codes,'' arXiv:2504.20847 (2025). \href{https://doi.org/10.48550/arXiv.2504.20847}{doi:10.48550/arXiv.2504.20847}; \href{https://arxiv.org/abs/2504.20847}{arXiv:2504.20847}.

  \item[\textup{[AMH26]}]\label{ref:ddfb-amh}
  G. Angl\`es Munn\'e and F. Huber, ``SDP bounds on quantum codes: rational certificates,'' arXiv:2603.19901 (2026). \href{https://doi.org/10.48550/arXiv.2603.19901}{doi:10.48550/arXiv.2603.19901}; \href{https://arxiv.org/abs/2603.19901}{arXiv:2603.19901}.

  \item[\textup{[HLZ25]}]\label{ref:ddfb-hlz}
  X. He, S. Lu, and B. Zeng, ``Co-Designing Quantum Codes with Transversal Diagonal Gates via Multi-Agent Systems,'' arXiv:2510.20728 (2025), version~3 (2026). \href{https://doi.org/10.48550/arXiv.2510.20728}{doi:10.48550/arXiv.2510.20728}; \href{https://arxiv.org/abs/2510.20728v3}{arXiv:2510.20728v3}.
\end{enumerate}

\subsection*{Comment}

Literature checked through 15 September 2026: no exact six-qubit code
satisfying Eqs.~\eqref{eq:ddfb-kl} and \eqref{eq:ddfb-transversal} and no
general exclusion was located.  This is a finite-size feasibility problem
derived from the unfilled $\mathrm C_{16}$ entry of Table~2 in
\sourcecite{ref:ddfb-zwhz}{ZWHZ25}, not a named nonexistence conjecture.  The
seven-qubit constructions, the restricted seven-qubit classification, the
distance-two catalogues, and the numerical six-qubit optimizations do not
decide it.  A negative answer must cover arbitrary local factors $U_j$ and every
$((6,2,3))_2$ code subspace, each of which is necessarily impure; a positive
answer must also certify genuine non-additivity, for example through a
noninteger value of Eq.~\eqref{eq:ddfb-signature}.

\subsection*{Field}

Quantum Error Correction

\subsection*{Topic}

Quantum coding theory

\subsection*{ID}

\texttt{op\_ddfbbe7ca0ead0c8}
